\subsection{综合评价的排序可信度}
问题五的综合评价只用于三种情景之间的相对比较。四个指标经正向化后，熵权、CRITIC、等权和 2000 次局部扰动都给出同一排序，其中“禁渔+污染+入侵”始终最差，复合压力的负向作用在当前比较框架下较为稳定。

综合分数不能替代单指标解释。食源、江豚和长江鲟仍对应不同层面的生态状态；若后续补入独立生态监测量并扩展污染、入侵强度梯度，排序结论的外推边界会更清楚。

\noindent 五个子问题共享同一条“恢复收益--上传导--长期复杂化--复合压力检验”的逻辑链。问题一和问题二给出收益出现的位置及其削弱机制，问题三和问题四刻画承载力提高后的行为变化方向，问题五则检验这部分收益在污染与入侵叠加下还能保留多少。我们统一时间口径、情景设定和上游输入，使不同模块之间能够互相传递信息而不至于脱节。

\noindent 按现有资料，模型目前更适用于方向性政策比较与机制解释。若后续补入连续河段监测、标准化 CPUE、珍稀物种年度调查以及污染分级资料，单区框架可以继续推进到分段比较，长期动力学部分也能与实际营养级数据建立更稳定的对应关系。

